% !TeX program = xelatex
% 编译方式：XeLaTeX，建议连续编译两次以更新总页数。
% 本文件不依赖外部图片、字体、参考文献或其他资源。
\documentclass[UTF8,a4paper,11pt,fontset=fandol]{ctexart}

\usepackage[left=20mm,right=20mm,top=19mm,bottom=20mm]{geometry}
\usepackage[table]{xcolor}
\usepackage{amsmath,amssymb}
\usepackage{booktabs,longtable,tabularx,array}
\usepackage{enumitem}
\usepackage{fancyhdr}
\usepackage{lastpage}
\usepackage{hyperref}
\usepackage{fvextra}
\usepackage{microtype}
\usepackage{titlesec}

\definecolor{NKBlue}{HTML}{17365D}
\definecolor{NKBlueTwo}{HTML}{244062}
\definecolor{NKLight}{HTML}{D9E2F3}
\definecolor{NKGray}{HTML}{666666}

\hypersetup{
  unicode=true,
  colorlinks=true,
  linkcolor=NKBlue,
  urlcolor=NKBlueTwo,
  pdfauthor={NKAI},
  pdftitle={第六届计图人工智能挑战赛赛题二代码提交说明文档},
  pdfsubject={Jittor原生点云降噪方法与复现说明},
  pdfkeywords={Jittor, 点云降噪, NKAI}
}

\setlength{\parindent}{2em}
\setlength{\parskip}{0.35em}
\setlist[itemize]{leftmargin=2.2em,itemsep=0.2em,topsep=0.25em}
\setlist[enumerate]{leftmargin=2.4em,itemsep=0.2em,topsep=0.25em}
\renewcommand{\arraystretch}{1.22}
\emergencystretch=2em

\titleformat{\section}
  {\Large\bfseries\color{NKBlue}}
  {}{0pt}{}
\titleformat{\subsection}
  {\large\bfseries\color{NKBlueTwo}}
  {}{0pt}{}
\titlespacing*{\section}{0pt}{1.2em}{0.55em}
\titlespacing*{\subsection}{0pt}{0.9em}{0.4em}

\pagestyle{fancy}
\fancyhf{}
\fancyfoot[C]{\small\color{NKGray}第 \thepage\ 页 / 共 \pageref{LastPage} 页}
\renewcommand{\headrulewidth}{0pt}
\renewcommand{\footrulewidth}{0pt}
\setlength{\headheight}{15pt}
\fancypagestyle{plain}{
  \fancyhf{}
  \fancyfoot[C]{\small\color{NKGray}第 \thepage\ 页 / 共 \pageref{LastPage} 页}
  \renewcommand{\headrulewidth}{0pt}
  \renewcommand{\footrulewidth}{0pt}
}

\DefineVerbatimEnvironment{CodeBlock}{Verbatim}{
  fontsize=\small,
  breaklines=true,
  breakanywhere=true,
  frame=single,
  framerule=0.35mm,
  framesep=2.5mm,
  rulecolor=\color{NKLight},
  formatcom=\color{black}
}

\newcommand{\code}[1]{\nolinkurl{#1}}

\begin{document}

\begin{center}
  {\zihao{2}\bfseries 第六届计图人工智能挑战赛}\par
  \vspace{0.35em}
  {\zihao{2}\bfseries 赛题二点云降噪方法与代码说明}\par
  \vspace{1.2em}
  {\zihao{4}\color{NKBlue}队伍：NKAI \quad 排名：031 \quad A榜成绩：79.23}\par
  \vspace{0.8em}
  \rule{0.88\textwidth}{0.6pt}
\end{center}

\section{一、队伍信息}

\begin{tabularx}{\textwidth}{>{\bfseries}p{3.2cm}X}
队伍名称 & NKAI \\
A榜排名 & 31（提交文件名使用三位字段 \code{031}） \\
A榜成绩 & \textbf{79.23 / 67.42 / 91.04}（Total / CD / P2S） \\
主要联系人 & 李佳璞，电话/微信：18883873159 \\
团队成员 & 刘迪乘，电话/微信：18570499315 \\
团队成员 & 曲恒睿，电话/微信：17606383128 \\
提交文件名 & \code{contest2_NKAI_031.zip} \\
\end{tabularx}

\section{二、我们的方法}

我们把赛题中的降噪任务看作带索引约束的点位移估计：输入和输出都保留 50,000 个点，输出中的第 $i$ 个点只对应输入中的第 $i$ 个点。这样既能修正局部噪声，又不会在拼接过程中改变点数或打乱顺序。

我们的主体网络由四个连续的去噪模块组成。每个模块先用动态边卷积提取局部邻域特征，再用状态空间模块补充长距离信息，最后预测该阶段的点位移。后一模块接收前一模块修正后的点云，因此网络可以逐步处理从明显离群点到细小表面抖动的不同误差。

状态空间部分由我们直接用 Jittor 实现。选择性扫描的前向和反向均通过 \code{jt.code} 编写 CUDA 算子；训练和推理不调用 PyTorch、Triton、\code{mamba_ssm} 或 \code{causal_conv1d}。为了控制显存，我们每隔 16 个序列位置保存一次状态，反向计算时从这些状态继续。提交模型使用 16 个状态维度和 float32。

单一模型很难同时兼顾桌类的大平面、沙发类的软边界以及低样本类别。我们因此从同一个基础模型出发，按数据形态训练密度分支和类别专家。各分支共享网络结构，只改变训练子集、噪声范围和密度感知损失权重。推理时，我们再用一个只依赖含噪点云的参考分支判断专家位移是否过大。这样做不需要测试集真值，也不会读取测试集 clean 点云或 mesh。

参考分支采用 Jittor 原生 StraightPCF。我们先训练速度模块，再训练级联速度模块和完整点云滤波模块。测试时执行两次滤波，并根据第一遍输出的均值、方差和局部变异系数确定第二遍强度。这个分支不负责替代主体网络，而是给专家输出提供尺度参照和异常回退。

\newpage
\section{三、数据准备与训练}

\subsection{从 mesh 构建训练点云}

我们的训练输入来自原始 ShapeNet mesh：

\begin{CodeBlock}
<TRAIN_ROOT>/shapenet/<synset>/<model>/models/model_normalized.obj
\end{CodeBlock}

我们在每个三角面上按面积分配采样概率，确定性地采样 50,000 个表面点。数据目录若带有 \code{model_normalized.obj.cache.npz}，我们读取其中的顶点和面，以保持 mesh 解析顺序；没有缓存时，我们使用代码包中的标准库 OBJ 解析器。采样随机数由公开 key 派生，因此重复运行会得到同一份 clean 点云。

训练时，我们在线加入 Gaussian 或 Laplace 噪声、2\% 离群点以及尺度变化，再从 50K 点云中抽取局部 patch。基础训练覆盖 15,733 个公开训练 mesh；类别训练使用包内固定 key list，列表只保存 \code{shapenet/<synset>/<model>} 标识，不包含坐标、标签、预测或权重。

\subsection{主体网络训练}

我们从随机初始化训练基础模型 100 个 epoch。单卡 micro-batch 为 4，累计 8 次梯度后更新一次参数，等效 batch 为 32。后续分支使用 micro-batch 4、累计 4 次梯度，等效 batch 为 16。共同参数为 \code{patch_size=1000}、\code{patch_ratio=1.2}、四个去噪模块、\code{frame_knn=32} 和梯度裁剪 1.0。

下表按用途列出流水线中的训练任务。表中的“样本数”表示 mesh 数乘以每个 mesh 的采样次数，$\lambda_{\mathrm d}$ 是密度感知 Chamfer 项的权重。

\begin{longtable}{>{\raggedright\arraybackslash}p{3.5cm}p{1.5cm}p{2.4cm}p{2.2cm}p{2cm}p{2.2cm}}
\toprule
\textbf{用途} & \textbf{轮数} & \textbf{样本数} & \textbf{学习率} & $\boldsymbol{\lambda_{\mathrm d}}$ & \textbf{随机种子} \\
\midrule
\endhead
基础模型 & 100 & $15733\times4$ & $10^{-4}$ & 0 & 2020 \\
类别均衡训练 & 3 & $650\times4$ & $10^{-5}$ & 0 & 20260814 \\
密度分支 & 2 & $3000\times1$ & $5\times10^{-6}$ & 0.5 & 20261303 \\
桌类初始分支 & 2 & $1200\times2$ & $5\times10^{-6}$ & 0 & 20260903 \\
桌类密度训练 & 2 & $1200\times1$ & $2\times10^{-6}$ & 0.5 & 20261411 \\
桌类强去噪训练 & 2 & $1200\times1$ & $10^{-6}$ & 1.0 & 20261511 \\
沙发强去噪训练 & 2 & $1000\times1$ & $2\times10^{-6}$ & 1.0 & 20262011 \\
飞机专家 & 2 & $1000\times4$ & $10^{-6}$ & 0.5 & 20262701 \\
桌类专家 & 2 & $1200\times4$ & $10^{-6}$ & 1.0 & 20262702 \\
沙发专家 & 2 & $1000\times4$ & $10^{-6}$ & 1.0 & 20262703 \\
低样本类别专家 & 2 & $1885\times4$ & $10^{-6}$ & 0.5 & 20262731 \\
\bottomrule
\end{longtable}

飞机、桌类和沙发专家使用纯 Laplace 噪声。三个区间分别为 $[0.0124896924,0.0146169608]$、$[0.0090486185,0.0129036043]$ 和 $[0.0083744542,0.0113465701]$。七个低样本类别各自使用独立区间，具体数值已经写入 \code{configs/} 下的类别噪声配置，流水线会自动读取。

我们的训练关系可以概括为：

\begin{CodeBlock}
基础模型
`-- 类别均衡训练
    |-- 密度分支
    |   |-- 沙发强去噪训练 -- 沙发专家
    |   |-- 飞机专家
    |   `-- 低样本类别专家
    `-- 桌类初始分支 -- 桌类密度训练
        `-- 桌类强去噪训练 -- 桌类专家
\end{CodeBlock}

\subsection{参考分支训练}

StraightPCF 的三个训练阶段依次为速度模块、级联速度模块和完整滤波模块，默认轮数为 3、30、40。每个 epoch 使用 10,000 个样本，batch size 为 32，patch size 为 1,000，Adam 学习率为 $10^{-4}$。我们从 mesh 采样 32,768 个表面点，并保留最多 1,024 个原始顶点；训练噪声为 $[0.008,0.014]$ 的 Laplace 分布，离群点比例为 2\%，离群尺度为 3。

\section{四、推理和结果融合}

\subsection{50K 点云推理}

测试数据目录为：

\begin{CodeBlock}
<TEST_ROOT>/shapenet/<synset>/<model>/noisy.npy
\end{CodeBlock}

我们先在 50K 点云上选择 150 个最远点采样中心，再为每个中心建立 2,000 点 patch。主体网络分别处理这些 patch；同一个原始点若出现在多个 patch 中，我们选择归一化中心距离最小的预测，因此每个输入索引只保留一个位移。该过程不会重新采样输出点。

参考分支执行两遍 StraightPCF。第一遍估计噪声尺度，第二遍按尺度调整滤波强度，随后以 $\gamma=0.5$ 融合两遍残差。参考输出与主体输出都保持输入索引。

\subsection{仅含噪输入的位移约束}

设含噪点为 $\mathbf{x}^{\mathrm n}_i$，参考分支输出为 $\mathbf{x}^{\mathrm r}_i$，类别专家输出为 $\mathbf{x}^{\mathrm e}_i$。我们计算

\[
r_i=\frac{\lVert\mathbf{x}^{\mathrm e}_i-\mathbf{x}^{\mathrm n}_i\rVert_2}
{\max\!\left(\lVert\mathbf{x}^{\mathrm r}_i-\mathbf{x}^{\mathrm n}_i\rVert_2,\,\mathrm{tiny}_{64}\right)},
\qquad q=Q_{0.95}(\{r_i\}).
\]

当 $q\le 4.0$ 时，我们采用

\[
\mathbf{x}^{\mathrm out}=\operatorname{float32}\!\left(
0.25\,\mathbf{x}^{\mathrm r}+0.75\,\mathbf{x}^{\mathrm e}
\right);
\]

否则使用参考分支输出。这个判断只比较三份由 noisy 输入得到的点云，不使用 clean 点云、标签或官方指标。

\subsection{按类别使用模型}

我们把提交时的模型选择直接写成固定类别映射，没有在测试集上搜索阈值。飞机、沙发和七个低样本类别使用相应专家与上面的位移约束；五个对密度更敏感的类别使用密度分支；其余类别使用类别均衡模型。

桌类同时保留桌类强去噪输出 $\mathbf{x}^{\mathrm b}$ 和经过位移约束的桌类专家输出 $\mathbf{x}^{\mathrm t}$。我们使用

\[
\mathbf{x}^{\mathrm table}=\operatorname{float32}\!\left(
\mathbf{x}^{\mathrm b}+1.25(\mathbf{x}^{\mathrm t}-\mathbf{x}^{\mathrm b})
\right).
\]

全部类别列表都随代码提交。我们在 \code{assemble.py} 中按 key 所属类别读取相应列表并完成组合，不使用测试点云以外的信息。

\subsection{结果格式}

我们把每个结果保存为 finite \code{float32 (50000,3)}：

\begin{CodeBlock}
shapenet/<synset>/<model>/denoised.npy
\end{CodeBlock}

打包时，我们按 key 排序并检查成员唯一性、路径安全、数组形状、数据类型和有限值。输出 ZIP 不包含目录穿越项或多余文件。

\section{五、环境与代码结构}

\subsection{运行环境}

我们使用的检查环境如下：

\begin{itemize}
  \item Ubuntu 22.04；
  \item NVIDIA RTX 4090 24GB，compute capability 8.9；
  \item CUDA Toolkit 12.4，cuDNN 8.9.7；
  \item Python 3.10.14；
  \item Jittor 1.3.11.0，NumPy 1.24.4。
\end{itemize}

推荐直接创建随包环境：

\begin{CodeBlock}
conda env create -f environment.yaml
conda activate <ENVIRONMENT_NAME>
\end{CodeBlock}

若机器已经安装兼容的 CUDA 和 cuDNN，也可以只安装正式运行依赖：

\begin{CodeBlock}
pip install -r requirements.txt
\end{CodeBlock}

我们把依赖分成两层：\code{requirements.txt} 只列 Jittor 和 NumPy；\code{environment.yaml} 另外固定 CUDA Toolkit、cuDNN 及赛事检查环境中的常用科学计算包。正式训练和推理只导入 Jittor 与 NumPy，PyTorch、Triton、SciPy、trimesh、OmegaConf 和 tqdm 均不在运行路径中。

Conda CUDA 把 \code{libcudart.so} 放在环境的 \code{lib/} 中，而 Jittor 1.3.11 默认检查 \code{lib64/}。我们在 \code{main.py} 中检测当前 Conda 环境，并在需要时建立 \code{lib64 -> lib} 兼容链接。

\subsection{提交包}

\begin{CodeBlock}
contest2_NKAI_031.zip
|-- code/
|   |-- main.py
|   |-- train.py
|   |-- train_rot.py
|   |-- infer.py
|   |-- generate_rot.py
|   |-- assemble.py
|   |-- package_results.py
|   |-- rot_jittor/
|   |-- configs/
|   |-- tools/
|   `-- tests/
|-- requirements.txt
|-- environment.yaml
`-- 提交说明文档.pdf
\end{CodeBlock}

我们没有在代码包中放入权重、训练点云、测试预测、mesh、日志或缓存。所有 checkpoint 都在指定工作目录中从随机初始化产生。

\section{六、复现步骤}

检查人员可以自行选择数据和输出目录。进入 \code{code/} 后，将下面三个变量改成机器上的实际绝对路径；其中 \code{NKAI_WORK_ROOT} 需要可写，并留出训练数据、checkpoint 和预测结果所需空间。

\begin{CodeBlock}
cd code
export NKAI_TRAIN_ROOT="/path/to/track2_train"
export NKAI_TEST_ROOT="/path/to/track2_test"
export NKAI_WORK_ROOT="/path/to/work_directory"
\end{CodeBlock}

先检查 Python、Jittor、CUDA 和自定义算子：

\begin{CodeBlock}
python main.py doctor
python main.py doctor --deep
\end{CodeBlock}

查看流水线将要执行的命令：

\begin{CodeBlock}
python main.py pipeline \
  --config configs/reproduce_full.json \
  --log-dir "$NKAI_WORK_ROOT/logs" \
  --dry-run
\end{CodeBlock}

确认数据路径后执行训练、推理、组合和打包：

\begin{CodeBlock}
python main.py pipeline \
  --config configs/reproduce_full.json \
  --log-dir "$NKAI_WORK_ROOT/logs"
\end{CodeBlock}

我们为每个耗时步骤指定了完成 checkpoint 或 manifest。重复执行同一命令时，流水线会跳过已经完成的步骤；如果需要重跑某一步，删除该步对应的完成文件即可。StraightPCF 训练支持从最近的训练状态继续。

若只需重新打包已经生成的预测，可以运行：

\begin{CodeBlock}
python main.py package-results -- \
  --root <PREDICTION_ROOT> \
  --key-list <TEST_KEY_LIST> \
  --zip <RESULT_ZIP> \
  --audit <AUDIT_JSON> \
  --expected-count 200 \
  --expected-points 50000
\end{CodeBlock}

\newpage
\section{七、运行验证}

我们在 Ubuntu 22.04、RTX 4090、CUDA 12.4.99、cuDNN 8.9.7、Python 3.10.14 和 Jittor 1.3.11.0 环境中解压代码并完成了以下检查：

\begin{itemize}
  \item 编译全部 Python 文件，并运行源码策略和结果融合单元测试；
  \item 编译选择性扫描 CUDA 算子，执行前向、反向和梯度检查；
  \item 从随机初始化更新主体网络参数，并检查密度感知损失和梯度累计；
  \item 依次训练 StraightPCF 的三个模块，保存、载入并继续训练；
  \item 对真实 50K 点云运行主体网络和两遍参考滤波；
  \item 组合单个测试 key，并生成带 CRC 和成员清单的结果 ZIP；
  \item 展开整条流水线，核对数据、训练、推理和打包命令之间的路径。
\end{itemize}

主体网络的 50K 实测输出为 finite float32，形状为 $(50000,3)$，没有触发 patch 回退；参考分支输出同样满足点数、类型和有限值要求。我们的源码审计还会检查禁止后缀、私有路径、密钥模式和禁止运行时依赖。

不同 GPU 上的动态 KNN 临界近邻和 CUDA reduction 次序可能产生很小的浮点差异，因此重新训练所得权重不要求逐字节相同。我们固定了训练 key、随机种子、噪声区间、网络结构、优化参数、类别映射和结果格式；提交结果仍需通过 shape、dtype、index、finite 和 ZIP 成员检查。

\section{八、来源与许可证}

我们参考了 3DMambaIPF、Mamba、IterativePFN 和 StraightPCF 的公开论文与实现，并使用 Jittor 完成正式训练和推理。来源说明、许可证文本和我们重写的模块边界见 \code{code/LICENSES.md} 与 \code{code/NOTICE}。

\vfill
\begin{center}
  \color{NKGray}\small
  NKAI · 第六届计图人工智能挑战赛赛题二
\end{center}

\end{document}
